Wir nehmen für den Radius des Rades $r = \raO$ an. Rollt das Rad ohne zu rutschen, so bewegt sich ein Punkt der Lauffläche gegenüber der Achse mit der Fahrgeschwindigkeit $v = \vaO$ im Kreis.
\begin{abcliste}
\abc Während einer Umdrehung fährt das Auto um den Radumfang weiter:
\begin{align}
 T &= \TaF \\
   &= \frac{2\pi\cdot\ra}{\va} \\
   &= \Ta \approx \result{\TaP{-3}{2}}
\end{align}

\abc Die Radialbeschleunigung eines Pneustücks (gegenüber der Achse) ist
\begin{align}
 a_\text{r} &= \apnF \\
   &= \frac{\left(\va\right)^2}{\ra} \\
   &= \apn \approx \result{\apnP{0}{2}}
\end{align}
Das ist etwa $\apgP{0}{2}$-mal so viel wie die Fallbeschleunigung $g$.
\end{abcliste}
Mit einem anderen vernünftigen Radius erhält man ähnliche Werte; es gilt $T \propto r$ und $a_\text{r} \propto 1/r$.
